\begin{titlepage}
\begin{center}

\begin{figure}[t]
\centering\includegraphics[width=0.15\textwidth]{singidunum-logo.png}
\end{figure}

{\Huge \textbf{\textsc{Univerzitet Singidunum}}}\\[0.3cm]
{\Large \textbf{Departman za poslediplomske studije}}

\vspace{2cm}

{\huge \textbf{\textsc{Doktorske studije}}}\\[0.2cm]
{\Large \textbf{Napredni sistemi zaštite}}

\vfill
{\Huge \textbf{Dizajn distribuirane platforme za e-trgovinu na javnoj \textit{blockchain} mreži zasnovan na optimizaciji računarskog napora radi smanjenja troškova naknada korisnika}}
\\[0.5cm]

{\Large \textbf{- doktorska disertacija -}}

\vfill

\begin{minipage}{0.45\textwidth}
    \raggedright
    \textbf{Mentor:}\\
    prof. dr Mladen Veinović
\end{minipage}
\hfill
\begin{minipage}{0.45\textwidth}
    \raggedleft
    \textbf{Student:}\\
    Vukašin S. Tasić
\end{minipage}

\vspace{1.5cm}
Beograd, 2026. godina

\end{center}
\end{titlepage}


%\cleardoublepage

\section*{Sažetak}
TODO

\section*{Abstract}
TODO

\newpage

\section*{Predgovor}
U uvodnoj glavi uspostavlja se konceptualni okvir za analizu savremenih sistema elektronske trgovine. Razmatraju se radni i informacioni sistemi, kao i formalna izražajnost poslovnih zahteva koje takvi sistemi treba da podrže. Dalje se analiziraju osnovni arhitektonski obrasci organizacije informacionih sistema, uz naglasak da većina savremenih rešenja zadržava centralizovan model upravljanja, nezavisno od toga da li su implementirana na jednom ili više računarskih čvorova, kao i ograničenja takvog pristupa. U tom kontekstu uvodi se pojam decentralizacije kroz različite dimenzije, kao ključni teorijski okvir za razumevanje potencijala unapređenja autonomije i korisničkog iskustva. Ova osnova služi kao polazište za dalju analizu izgradnje platforme za e-trgovinu na javnim blokčejn tehnologijama, sa ciljem povećanja otpornosti i distribuirane kontrole, čime se uvodi problem transakcionih troškova za korisnike i definiše centralni istraživački fokus ovog doktorskog rada.

Druga glava uspostavlja analitički okvir za evaluaciju informacionih sistema kroz njihova tehnička svojstva i kompromis između performansi, skalabilnosti, dostupnosti i bezbednosti. Analiza obuhvata fundamentalne principe i ograničenja distribuiranih sistema, uključujući CAP teoremu, skalabilnost, zagušenje, režijsko opterećenje, koordinaciju, propagaciju, replikaciju, konsenzus i kontrolu konkurentnog pristupa, kao i njihove posledice po performanse i operativnu stabilnost sistema. Posebna pažnja posvećena je operativnim svojstvima informacionih sistema kroz bezbednosne modele, kriptografske mehanizme, autentikaciju, autorizaciju i evidenciju, čime se pokazuje da tehnička arhitektura sistema istovremeno određuje i njegove mogućnosti i njegova ograničenja.

\tableofcontents
\hypersetup{pageanchor=true}